\section{Conventional Prediction Methods} \label{section:conventional_methods}

% ---------- %

Understanding and predicting the behaviour of material interfaces requires a firm foundation in computational methods
that are rooted in quantum and statistical mechanics.

This section introduces the core approaches historically used to model interfaces at the atomic and mesoscale level,
focusing primarily on first-principles methods and classical simulation techniques.

Although each of these methods offers powerful insights, their accuracy, computational cost, and suitability for
interface modelling vary substantially.

% ---------- %

A primary tool in this field is \textit{Density Functional Theory} (DFT), a quantum-mechanical framework capable of
describing electronic structure with high precision.

DFT is well-suited to calculating key interfacial properties such as formation energy, charge redistribution, and band
alignment, making it an essential component of most modern interface studies.

Within this domain, software packages such as \textsc{VASP} have become widely adopted for their efficient
implementations of the Kohn-Sham formalism.

% ---------- %

However, DFT has limits. Computationally, DFT's cubic scaling with system size makes it computationally prohibitive for
large or disordered interface models.

As a result, researchers also turn to \textit{classical methods}, such as \textit{Molecular Dynamics} (MD).

MD enables the simulation of dynamic processes and thermal fluctuations over nanosecond timescales.

It is also adept at capturing thermal effects and atomic rearrangements.

% ---------- %

These conventional approaches provide the groundwork for predicting critical interfacial properties, including adhesion
energy, charge transfer behaviour, and electron or phonon transport.

They remain indispensable despite their limitations, particularly when used in tandem.

This section surveys these methods in turn, addressing their theoretical underpinnings, practical implementations, and
their relative advantages and constraints in the context of interface science.

% ---------- %

A foundational step in many interface modelling workflows is lattice matching between two crystalline materials.

This involves aligning compatible Miller planes to minimise strain and identify viable supercell geometries.

The method introduced by Zur et al.~(1984) laid the groundwork for many subsequent algorithms used in tools such as
\textsc{METADISE} and \textsc{ARTEMIS}, which systematically enumerate low-strain combinations, slips, terminations, and
atomic registries.

These approaches assume idealised surfaces and often neglect more complex processes such as surface reconstructions,
chemical interdiffusion, or local disorder.

% todo: Describe it (lattice matching / geometry-driven interface generation)
% don't just reference.

% ---------- %

DFT remains the most accurate and widely used quantum mechanical method for predicting interfacial properties.

It solves the Kohn--Sham equations to determine the electronic ground state, providing access to total energies, charge
densities, and band edge positions.

% ---------- %

MD simulates the atomic trajectories under classical Newtonian mechanics, governed by interatomic potentials.

Classical MD is well-suited to studying thermal processes such as diffusion, defect migration, and strain relaxation on
picosecond to nanosecond timescales.

% ---------- %

\subsection{Density Functional Theory (DFT)}

% ---------- %

Density Functional Theory (DFT) is a first-principles quantum mechanical method used extensively to model the electronic
structure of condensed matter systems, including molecules, solids, and, critically, material interfaces.

It provides a variational framework in which the ground-state properties of a many-electron system can be obtained from
a functional of the electron density $\rho(\mathbf{r})$, rather than from the many-body wavefunction
$\Psi(\mathbf{r}_1, \dots, \mathbf{r}_N)$.

This simplification, formalised through the Hohenberg-Kohn theorems, underpins the modern implementation of DFT via the
Kohn-Sham approach, where the interacting system is replaced by an auxiliary system of non-interacting electrons subject
to an effective potential that captures Coulomb, exchange, and correlation interactions.

% ---------- %

The total Kohn-Sham energy is given by:

\begin{equation} E_{\mathrm{KS}}[\rho] = T_s[\rho] + \int V_{\mathrm{ext}}(\mathbf{r}) \rho(\mathbf{r}) \, \mathrm{d}\mathbf{r} + \frac{1}{2} \int \int \frac{\rho(\mathbf{r}) \rho(\mathbf{r'})}{|\mathbf{r} - \mathbf{r'}|} \, \mathrm{d}\mathbf{r} \, \mathrm{d}\mathbf{r'} + E_{\mathrm{xc}}[\rho], \end{equation} where $T_s[\rho]$ is the kinetic energy of the non-interacting electrons, $V_{\mathrm{ext}}$ is the external (ionic) potential, and $E_{\mathrm{xc}}[\rho]$ is the exchange-correlation functional, the exact form of which remains unknown.

% ---------- %

Common approximations include the Local Density Approximation (LDA), Generalised Gradient Approximation (GGA), and
hybrid functionals.

While LDA assumes a locally uniform electron gas and often leads to overbinding, GGA incorporates density gradients and
improves accuracy for molecular and inhomogeneous systems.

Hybrid functionals, such as HSE06, incorporate a portion of exact exchange and offer improved bandgap and localisation
predictions at significantly greater computational cost.

% ---------- %

DFT has become foundational in materials modelling due to its ability to predict not only total energies and optimised
geometries, but also derived properties including electronic band structures, charge densities, dielectric responses,
phonon spectra, and interfacial energetics.

Its applications to interface modelling, however, are computationally constrained.

% ---------- %

All DFT simulations in this project are conducted using the Vienna Ab initio Simulation Package (VASP), a widely used
plane-wave-based implementation of Kohn-Sham DFT.

VASP employs the projector-augmented wave (PAW) method to model core-valence interactions, and supports a variety of
exchange-correlation functionals including GGA-PBE and hybrid HSE06.

Structural relaxations are typically performed with conjugate gradient algorithms until residual forces fall below a
chosen threshold (e.g. 0.01 eV/\AA).

% ---------- %

Periodic boundary conditions are employed in all three spatial directions, necessitating the introduction of vacuum
regions for slab models and sufficient supercell dimensions to minimise image-image interactions in interface
simulations.

% ---------- %

The simulation input is controlled via the INCAR, POSCAR, KPOINTS, and POTCAR files, with convergence typically assessed
by k-point density and energy cutoff tests. Pulay and Broyden mixing schemes are used for charge density convergence,
and dipole corrections may be employed for asymmetric slab geometries.

% ---------- %

While VASP provides accurate and robust results for interfacial modelling, it remains computationally intensive.

For extended or complex interfacial systems, this motivates the introduction of machine-learned potentials (MLPs), which
aim to reproduce DFT accuracy.

% ---------- %

\subsection{Other Modelling Methods}

% ---------- %

In addition to DFT, other classes of atomistic modelling, such as MD, has proven valuable for exploring material
interfaces, particularly when temporal evolution or thermal effects are of interest.

% ---------- %

Molecular Dynamics is a time-resolved simulation method in which the trajectories of atoms are computed by numerically
integrating Newton's equations of motion.

Forces on each atom are derived either from empirical interatomic potentials (classical MD) or quantum-mechanical
calculations (ab initio MD, e.g. Car-Parrinello or Born-Oppenheimer approaches).

In interface science, MD is used to study thermal stability, strain relaxation, diffusion, and defect evolution at
finite temperatures.

% ---------- %

One advantage of MD is its capacity to resolve dynamic processes at the atomic level, such as dislocation glide, grain
boundary motion, or interfacial melting.

However, timescales are severely limited: typical simulations span nanoseconds, far shorter than the seconds to minutes
often relevant in experiments.

Classical MD can scale to millions of atoms but sacrifices accuracy; ab initio MD is more accurate but typically limited
to fewer than 1,000 atoms and short timescales due to computational cost.

% ---------- %

For interface modelling, MD is particularly useful when studying thermally activated restructuring, interdiffusion, and
local vibrational behaviour, especially under conditions such as annealing or irradiation.

It also serves as a preliminary stage in workflows involving machine-learned potentials, providing dynamical data for
training.

% ---------- %

MD complements DFT in the modelling of interfaces.

While DFT offers high-fidelity energetics and electronic structure, MD enables atomistic dynamics over nanoseconds.

As such, hybrid approaches are increasingly common, with DFT used to parameterise MD models, or machine learning
techniques introduced to accelerate sampling across scales.

% ---------- %

\subsection{What are the advantages and limitations of DFT and MD?}

% ---------- %

Each computational method traditionally employed in the modelling of interfaces, DFT and MD, offers a distinct balance
between accuracy, computational efficiency, and physical scale.

While often applied in isolation, their complementary regimes suggest the need for hybrid or hierarchical approaches in
the study of real interfaces.

% ---------- %

\textbf{Density Functional Theory (DFT)} is the most widely used \textit{ab initio} method for predicting atomic-scale
properties at interfaces.

Crucially, DFT is capable of modelling materials across the periodic table and is the backbone of modern electronic
structure prediction.

Its limitations include conventional exchange-correlation functionals (e.g. GGA) systematically underestimate band gaps
and cannot capture excited-state dynamics without costly extensions such as GW or time-dependent DFT.

% ---------- %

\textbf{Molecular Dynamics (MD)} offers time-resolved access to thermally activated processes and interfacial dynamics.

Classical MD, using empirical or machine-learned interatomic potentials, enables simulation of thousands to millions of
atoms over nanosecond to microsecond timescales.

However, MD inherits limitations from its force field or potential model: empirical potentials may not transfer well
across bonding environments, while high-fidelity machine-learned potentials require extensive and chemically diverse
training data.

Furthermore, standard MD remains constrained to accessible timescales, making it ill-suited for capturing rare events
such as defect migration, reconstruction, or long-term phase evolution.

% ---------- %

\subsection{How are properties like adhesion energy or electron transport predicted?}

% ---------- %

\textbf{Adhesion energy} is a thermodynamic quantity describing the work required to separate two materials at their
interface.

It is commonly computed from the total energies of the isolated constituents and the relaxed interface configuration.

Within DFT, the adhesion energy $E_\mathrm{adh}$ is given by:

\[ E_\mathrm{adh} = \frac{1}{A} \left( E_\mathrm{slab1} + E_\mathrm{slab2} - E_\mathrm{interface} \right), \]

where $A$ is the interface area, and the energies refer to fully relaxed geometries of the individual slabs and the
interface supercell.

% ---------- %

% todo: add interface energy
% todo: add surface energy
% todo: add strain energy
% todo: add strain: tension, compression, shear
% todo: add formation energy

% ---------- %

\textbf{Electron transport} across interfaces is more nuanced and often requires evaluation of both band alignment and
carrier mobility.

Band alignment can be predicted via several approaches.

The simplest, Anderson's rule aligns vacuum levels to estimate offsets from electron affinities or ionisation
potentials.

However, this often fails to capture interface-specific effects such as charge transfer, built-in fields, or
hybridisation.

More accurate predictions therefore rely on self-consistent DFT calculations that extract the local electrostatic
potential across the heterostructure to determine band offsets directly.

% ---------- %

For example, the offset between conduction or valence band edges across an interface can be obtained by aligning the
macroscopic average electrostatic potentials of the constituent slabs with that of the full interface.

This method accounts for relaxation, dipole formation, and interfacial reconstruction.

In cases where electronic coupling is strong, or the valence states are spatially delocalised across layers (as in
type-II or type-III band alignment), transport characteristics such as charge separation, carrier confinement, or
tunnelling probability can change significantly.

% ---------- %

Charge transport properties such as electrical conductivity, Seebeck coefficient, or mobility may be
approximated from the band dispersion near the Fermi level using Boltzmann transport theory, or more rigorously via
non-equilibrium Green's function (NEGF) methods, especially for low-dimensional or nanoscale interfaces.

% ^^^ note: Intro